\documentclass[10pt,a4paper]{amsart}
\usepackage[margin=19mm]{geometry}
\usepackage{amsmath,amssymb,mathtools}
\usepackage[scr=rsfs,cal=euler]{mathalfa}
\usepackage{xfrac}
\usepackage[unicode,hidelinks,bookmarks=false]{hyperref}
\newcommand{\ie}{i.e.\ }
\newcommand{\cf}{cf.\ }
\newcommand{\resp}{respectively\ }
\newcommand{\Subsection}{\subsection}
\providecommand{\nobreakdash}{\nobreak\hbox{-}}
\setlength{\emergencystretch}{2em}
\multlinegap0pt
\newtheorem{prop}{Proposition}
\newtheorem{defin}{Definition}
\newcommand{\A}{\mathcal{A}}
\newcommand{\hp}{{\mathbb{P}^1(\mathbb{H})}}
\newcommand{\ordt}{\operatorname{ord}^t}
\newcommand{\s}{{\mathbb{S}}}
\title{Quaternionic Nevanlinna Functions}
\author{Muhammad Ammar}
\date{Source: arXiv:2603.19332v1 (CC BY 4.0)}
\begin{document}
\maketitle

\noindent\emph{Source and licence.} Quaternionic Nevanlinna Functions. Version arXiv:2603.19332v1 (CC BY 4.0), distributed by arXiv under the Creative Commons Attribution 4.0 International licence, http://creativecommons.org/licenses/by/4.0/, as recorded in the arXiv licence metadata for this exact version and shown on the abstract page https://arxiv.org/abs/2603.19332v1. Full licence notice: \texttt{licenses/arxiv-2603.19332-CC-BY-4.0.txt}.\par\smallskip

\noindent\textit{Editorial scope.} Counted unit: the article's prop prop (source0.tex lines 527--532). The article's complete original proof of it is delivered with it (source0.tex lines 533--559, one \texttt{proof} environment of 27 lines). No mathematics is written, repaired or re-derived: the statement and the proof are the article's own lines, in the article's own notation, and no retained line is substituted or re-flowed.  Retained uncounted as the source-local context the proof needs: lines 439--442; lines 469--476; lines 499--503.  Nothing was omitted from any retained span, so the package carries no omission record.  No retained line contains a citation, so the wrapper carries no bibliography and no bibliography entry.  No retained line contains a figure environment, an image-inclusion command or drawing code, so no image is delivered and the package declares no asset dependency.  Editorial formatting only, in addition: an AMS \texttt{amsart} preamble that restates the article's own declarations and macros used by the retained lines, namely the environments \texttt{prop}, \texttt{defin} on separate counters, together with the standard packages those lines need.  The retained environments print in this document as Proposition 1, Defin 1 in order of appearance, whereas the article prints the same items with the numbering of its own full document.  The complete native source of the article as distributed in the arXiv e-print for this exact version is bundled unchanged as \texttt{sources/196765/source0.tex}.
\begin{prop}\label{quatIntCou2}
    Let $f$ be defined as above. Then, \begin{align*}N(f,a,r)&=n(f,a,0)\log r + \int_0^r[n(f,a,t)-n(f,a,0)]
    \,\frac{dt}{t} \\&\quad+\sum_{\substack{ \s_\zeta\in\mathcal{S}(\overline{\mathbb B_r}) \\ \zeta\neq 0}}\ordt_{\s_\zeta}(f,a) \frac{|\zeta|^4-r^4}{4r^2|\zeta|^4}\left(2(\Re \zeta)^2-|\zeta|^2\right).\end{align*}
\end{prop}

\begin{prop}\label{quaIntCou3}
    Let $f$ be defined as in Proposition \ref{quatIntCou2}. Then, \begin{align*}
        N(f,a,r) &= n(f,a,0)\log r + \int_0^r[n(f,a,t)-n(f,a,0)]\,\frac{dt}{t} \\
        &\quad +  \int_0^r\frac{t^4+r^4}{2r^2t^3}[n(f,a,t)-n(f,a,0)]\,dt \\ 
        &\quad + \sum_{\substack{ \s_\zeta\in\mathcal{S}(\overline{\mathbb B_r}) \\ \zeta\neq 0}}
        \ordt_{\s_\zeta}(f,a) \frac{|\zeta|^4-r^4}{2r^2|\zeta|^4}(\Re \zeta)^2
    \end{align*}
\end{prop}

We are now left with a term that cannot be further simplified purely in terms of the radial unintegrated counting function, due to the factor of $\Re\zeta $. 
\begin{defin}[Angular Counting Term]\label{AngularCountingTerm}
     Let $f:\Omega_D\to\hp$ be semiregular, and let the closure of $\mathbb{B}_r$ be contained in $\Omega_D$. Then, we define the angular counting term as \[ \A(f,a,r)\coloneq \sum_{\substack{ \s_\zeta\in\mathcal{S}(\overline{\mathbb B_r}) \\ \zeta\neq 0}}
        \ordt_{\s_\zeta}(f,a) \frac{|\zeta|^4-r^4}{2r^2|\zeta|^4}(\Re \zeta)^2.\]
\end{defin}

\begin{prop}
    Let $f, \A$ be defined as in \ref{AngularCountingTerm}. Then, \begin{align*}
        \A(f,a,r)&=4r^2\iint_{0\leq h\leq t\leq r} ht^{-5}[a_t(f,a,h)-a_t(f,a,0)]\, dh\, dt \\&\quad - \int_0^r 2r^2t^{-3}[n(f,a,t)-n(f,a,0)]\,dt. \\
        &= -2r^2\int_0^rt^{-5}[a_r^\Re(f,a,t)-a_r^\Re(f,a,0)]\,dt.  \\
    \end{align*}
\end{prop}

\begin{proof}
    The proof is much the same as in Propositions \ref{quatIntCou2} and \ref{quaIntCou3}. We again begin by decomposing \begin{align}\label{quaIntCou4.1}\sum_{\substack{ \s_\zeta\in\mathcal{S}(\overline{\mathbb B_r}) \\ \zeta\neq 0}}\ordt_{\s_\zeta}(f,a) \frac{|\zeta|^4-r^4}{2r^2|\zeta|^4}(\Re \zeta)^2 
    &=\sum_{\substack{ \s_\zeta\in\mathcal{S}(\overline{\mathbb B_r}) \\ \zeta\neq 0}}\ordt_{\s_\zeta}(f,a)\frac{(\Re \zeta)^2}{2r^2} \nonumber\\
    &\quad- \sum_{\substack{ \s_\zeta\in\mathcal{S}(\overline{\mathbb B_r}) \\ \zeta\neq 0}}\ordt_{\s_\zeta}(f,a)\frac{(\Re \zeta)^2r^2}{2|\zeta|^4} \nonumber\\
    &= \sum_{\substack{ \s_\zeta\in\mathcal{S}(\overline{\mathbb B_r}) \\ \zeta\neq 0}}\ordt_{\s_\zeta}(f,a)\frac{(\Re \zeta)^2}{2r^2} \nonumber\\
    &\quad- \frac{r^2}{2}\sum_{\substack{ \s_\zeta\in\mathcal{S}(\overline{\mathbb B_r}) \\ \zeta\neq 0}}\ordt_{\s_\zeta}(f,a)\left(\frac{(\Re \zeta)^2}{r^4}+\int_{|\zeta|}^r\frac{4(\Re \zeta)^2}{t^5}\,dt\right) \nonumber\\
    &=-\frac{r^2}{2}\sum_{\substack{ \s_\zeta\in\mathcal{S}(\overline{\mathbb B_r}) \\ \zeta\neq 0}}\ordt_{\s_\zeta}(f,a)\int_{|\zeta|}^r\frac{4(\Re \zeta)^2}{t^5}\,dt, 
    \end{align}
so it suffices to look at the remaining term. We have \begin{align}\label{quaIntCout4.3}
    2r^2\sum_{\substack{ \s_\zeta\in\mathcal{S}(\overline{\mathbb B_r}) \\ \zeta\neq 0}}\ordt_{\s_\zeta}(f,a)\int_{|\zeta|}^r\frac{(\Re \zeta)^2}{t^5} &= 2r^2\int_0^r t^{-5}\left(\sum_{\substack{ \s_\zeta\in\mathcal{S}(\overline{\mathbb B_t}) \\ \zeta\neq 0}}\ordt_{\s_\zeta}(f,a)(\Re \zeta)^2\right) \, dt.
\end{align}
Observe that we can treat the inner term in a similar manner, in that we attempt to further write the term $(\Re \zeta)^2$ as an integral. We analyze this term independently, and thus we have \begin{align}\label{quaIntCou4.2}
\sum_{\substack{ \s_\zeta\in\mathcal{S}(\overline{\mathbb B_t}) \\ \zeta\neq 0}}\ordt_{\s_\zeta}(f,a)(\Re \zeta)^2 
&= \sum_{\substack{ \s_\zeta\in\mathcal{S}(\overline{\mathbb B_t}) \nonumber \\ \zeta\neq 0}}\ordt_{\s_\zeta}(f,a)\left(t^2 -\int_{|\Re \zeta|}^t2h \,dh\right) \\
&= t^2 \sum_{\substack{ \s_\zeta\in\mathcal{S}(\overline{\mathbb B_t}) \nonumber\\ \zeta\neq 0}}\ordt_{\s_\zeta}(f,a)-2\int_0^t h\left( \sum_{\substack{ \s_\zeta\in\mathcal{S}(\overline{\mathbb B_t}) \\ \zeta\neq 0 \\|\Re \zeta|\leq h}}\ordt_{\s_\zeta}(f,a)\right)\,dh \\ 
&=t^2[n(f,a,t)-n(f,a,0)]-2\int_0^t h [a_t(f,a,h)-a_t(f,a,0)]\,dh.
\end{align}
Thus, we have \begin{align*}
      2r^2\sum_{\substack{ \s_\zeta\in\mathcal{S}(\overline{\mathbb B_r}) \\ \zeta\neq 0}}\ordt_{\s_\zeta}(f,a)\int_{|\zeta|}^r\frac{(\Re \zeta)^2}{t^5} 
      &= 2r^2\int_0^rt^{-3}[n(f,a,t)-n(f,a,0)]\, dt \\
      &\quad- 4r^2\int_0^rt^{-5}\left(\int_0^th[a_t(f,a,h)-a_t(f,a,0)]\,dh \right)\,dt. 
\end{align*}
The final integral term on the right-hand side cannot be expressed as a single variable, because $a_t(f,a,h)$ depends on both $t$ and $h$, though we may still apply Fubini's theorem. The first equality follows by recalling equations \ref{quaIntCou4.1} and \ref{quaIntCou4.2}. 

The second equality comes from returning to Equation \ref{quaIntCout4.3}, where we recall the definition of $a_r^\Re (f,a,t)$ and note \[\sum_{\substack{ \s_\zeta\in\mathcal{S}(\overline{\mathbb B_t}) \\ \zeta\neq 0}}\ordt_{\s_\zeta}(f,a)(\Re \zeta)^2=a_r^\Re (f,a,t)-a_r^\Re (f,a,0).\] 

\end{proof}

\end{document}
